% ===============================================================
\chapter{Putnam, `Meaning and Reference'}
% ===============================================================

\section{Semantic internalism vs semantic externalism}

These are attitudes one could take towards meanings.

\begin{definition}\leavevmode
\begin{enumerate}[label=(\alph*)]
  \item \emph{Semantic internalism}: meanings are mental entities of some sort (thoughts, concepts).
  \item \emph{Semantic externalism}: meanings are mind-external entities of some sort (functions, sets, etc.).
\end{enumerate}
\end{definition}

\begin{example}\leavevmode
\begin{enumerate}[label=\roman*.]
  \item `Dog' viewed with semantic internalism would be the concept \textsc{dog}.
  \item `Dog' viewed with semantic externalism would be the kind \emph{canine}, or the function from possible worlds to the set of dogs.
\end{enumerate}
\end{example}

Frege opposed semantic internalism because communication would not be effective without a common meaning of the word `dog'. Semantic internalism is also known as \emph{psychologism} (what Frege was opposing).

\section{An argument for semantic internalism, and the Psychological Thesis}

\begin{argument}[For semantic internalism]\leavevmode
\begin{premises}
  \item[P1)] Sentences express thoughts.
  \item[P2)] Concepts are units, or constituents, of thoughts.
  \item[P3)] Meaning is compositional.
  \item[C1)] So, sentences couldn't express thoughts unless words expressed concepts.
  \item[C2)] Hence, words express concepts.
\end{premises}
\end{argument}

\begin{definition}[Psychological Thesis]
Meanings are psychological states, like thoughts.
\end{definition}

\section{Semantic externalism and the Determination Thesis}

\begin{enumerate}[label=(\alph*)]
  \item Functions relate each input to a unique output. Thus, identical functions with the same inputs must return identical outputs.
  \item For Frege, sense determines reference, so that the extension of an expression is a function of its sense.
\end{enumerate}

\begin{definition}[Determination Thesis]
Meaning determines semantic value. That is, expressions having the same meaning must also have the same semantic value, though expressions can have the same semantic value without having the same meaning.
\end{definition}

\section{Singular, general, natural kind and artifact terms}

\begin{definition}\leavevmode
\begin{enumerate}[label=(\alph*)]
  \item \emph{Singular terms}: expressions which refer (if at all) to one thing.
  \item \emph{General terms}: expressions which refer (if at all) to multiple things of the same kind; common nouns are good examples.
  \item \emph{Natural kind terms}: general terms referring to things which occur in nature and whose existence is not dependent on human nature.
  \item \emph{Artifact terms}: general terms referring to artifacts, or human creations.
\end{enumerate}
\end{definition}

\begin{note}
General terms can effectively act as a good replacement for names. For the most part, these are natural kind terms:
\begin{enumerate}
  \item We can existentially generalise over them (Water is on the floor $\Rightarrow$ something is on the floor).
  \item They are not scope ambiguous (I didn't see water on the floor).
\end{enumerate}
\end{note}

\section{The Twin Earth experiment}

The intention of this experiment is to show that meanings cannot be a function of anything psychological.

\begin{example}[Twin Earth]
Twin Earth is exactly like Earth at some point before the Scientific Revolution, and thus before the discovery of the chemical composition of water, say 1600. Earth and Twin Earth are exactly alike except in one respect: the stuff people on Earth call `water' is composed of H$_2$O, while the stuff people on Twin Earth call `water' is composed of something else, say XYZ. Despite this, people on Twin Earth have the exact same psychological states as people on Earth---they have the same beliefs, thoughts, concepts, etc. Yet `water' on Twin Earth has a different semantic value than `water' on Earth: on Earth, `water' describes quantities of H$_2$O, while on Twin Earth it describes quantities of XYZ.
\end{example}

Assumptions:
\begin{enumerate}
  \item Twin Earth is not a possible world, but another planet, which actually exists.
  \item Putnam assumes that psychological states are physically determined, so that people who are molecule-for-molecule identical have the same psychological states. Thus, there can be no difference in psychological states without a corresponding difference in neurophysiology.
\end{enumerate}

\section{Putnam's argument against semantic internalism}

\begin{argument}[Putnam]\leavevmode
\begin{premises}
  \item[P1)] Meaning determines semantic value.
  \item[C1)] Thus, if the semantic value of an expression in one place is different from its semantic value in another, then that expression must be ambiguous.
  \item[P2)] On Earth, the semantic value of `water' is H$_2$O. On Twin Earth, the semantic value of `water' is instead XYZ.
  \item[C2)] P2) $\Rightarrow$ `water' is ambiguous.
  \item[P3)] The psychological states and concepts of English speakers on Earth are the same as those of English speakers on Twin Earth.
  \item[C3)] Thus, speakers' psychological states or concepts do not determine semantic value.
  \item[C4)] Thus, meanings cannot be psychological states or concepts.
\end{premises}
\end{argument}

\section{The rigidity of natural kind terms}

\begin{enumerate}[label=(\alph*)]
  \item Kripke proposed that natural kind terms, such as `water', are rigid designators.
  \item Putnam further suggested that the meanings of natural kind terms have an indexical element to them.
  \item We can see the indexical in action in the following:
  \begin{enumerate}[label=\roman*.]
    \item In all possible worlds, the semantic value of `water' is whatever bears a similarity relation to what is water here, namely H$_2$O.
    \item We can use the similarity relation to create an equivalence class for water and H$_2$O.
  \end{enumerate}
\end{enumerate}

\section{Objections to Putnam}

\begin{objection}[Indexicals]\leavevmode
\begin{enumerate}[label=\roman*.]
  \item Putnam only shows that we must reject the \emph{combination} of the Determination Thesis and the Psychological Thesis.
  \item David Kaplan's indexical theory. According to Kaplan, there are three aspects of meaning:
  \begin{enumerate}[label=(\alph*)]
    \item \emph{Character}: together with facts about context, determines content.
    \item \emph{Content}: together with facts about the world, determines extension (Fregean sense).
    \item \emph{Extension}: semantic value.
  \end{enumerate}
  There is an aspect of meaning which does not determine semantic value, namely character. This threatens the Determination Thesis.
  \item The above implies that the Determination Thesis is false, and therefore meaning does not determine semantic value.
  \item In the argument against semantic internalism, we have falsified P1), which leaves us with P3) and C3):
  \begin{premises}
    \item[P3)] The psychological states and concepts of English speakers on Earth are the same as those of English speakers on Twin Earth.
    \item[C3)] So, psychological states or concepts do not determine semantic value.
  \end{premises}
  \item Even if C3) is true, the consequences of David Kaplan's theory still allow us to accept the Psychological Thesis:
  \begin{enumerate}
    \item Meaning is not content, but character.
    \item Thus, meaning determines content, not semantic value.
    \item Character consists in speakers' psychological states or concepts.
  \end{enumerate}
  \item This implies that Twin Earth is consistent with semantic internalism.
\end{enumerate}
\end{objection}

\begin{objection}[Psychological externalism]\leavevmode
\begin{definition}[Psychological externalism]
Psychological states and concepts are (at least partially) externally individuated. That is, the contents of beliefs and other psychological states, as well as the concepts which constitute them, are at least partially determined by one's physical environment.
\end{definition}
This has the capacity to cripple P3) in Putnam's argument against semantic internalism. Psychological externalism implies that:
\begin{enumerate}
  \item Suppose that while pointing at some H$_2$O, an Earth speaker forms the belief expressed by `This is water'. And suppose the same thing happens for a Twin Earth speaker.
  \item Given psychological externalism, these are different beliefs, as the concepts involved pick out different kinds of stuff: H$_2$O and XYZ.
  \item Thus, speakers on Earth and speakers on Twin Earth do not have the same psychological states, despite being molecular duplicates.
\end{enumerate}
The above follows from these two theses:
\begin{itemize}
  \item[(T1)] Psychological states are individuated by their contents. Thus, psychological states having different contents are distinct.
  \item[(T2)] Natural kind concepts are externally individuated. That is, the content of a concept like \textsc{water} is determined by one's physical environment.
\end{itemize}
By (T2), speakers on Earth possess different concepts than speakers on Twin Earth, because while \textsc{water}$_{e}$ picks out H$_2$O, \textsc{water}$_{te}$ picks out XYZ. So, by (T1), speakers' beliefs on Earth about water are different from those of speakers on Twin Earth. If so, then neurophysiology does not completely determine the contents of our psychological states.
\end{objection}
